% Mapa de dependencias de matematicado: una flecha X -> Y significa «Y necesita X».
\documentclass[tikz,border=4pt]{standalone}
\usepackage{resonancia}
\begin{document}
\begin{tikzpicture}[x=1cm, y=1cm, nodo/.style={caja, minimum width=2.5cm, minimum height=0.85cm, font=\sffamily\small}]
  % columnas por bloque (de izquierda a derecha, en orden de lectura)
  \node[nodo, draw=rGris, fill=rGrisClaro] (E)  at (0, 3.2) {E · Edificio};
  \node[cajaAzul, nodo] (A1) at (0, 1.6) {A1 · Matrices};
  \node[cajaAzul, nodo] (A2) at (3.3, 2.4) {A2 · Valores propios};
  \node[cajaAzul, nodo] (A3) at (3.3, 0.6) {A3 · SVD};
  \node[cajaVerde, nodo] (S1) at (6.6, 2.4) {S1 · Laplace};
  \node[cajaVerde, nodo] (S2) at (6.6, 0.6) {S2 · Estados};
  \node[cajaVerde, nodo] (S3) at (9.9, 1.5) {S3 · Transformada Z};
  \node[cajaNaranja, nodo] (F1) at (13.2, 2.4) {F1 · Fourier};
  \node[cajaNaranja, nodo] (F2) at (16.5, 3.2) {F2 · Hilbert};
  \node[nodo, draw=rVioleta, fill=rVioleta!10] (P1) at (16.5, 1.6) {P1 · Estocásticos};
  \node[nodo, draw=rVioleta, fill=rVioleta!10] (P2) at (16.5, 0.0) {P2 · Correlación};
  \foreach \a/\b in {E/A2, A1/A2, A1/A3, A2/A3, A2/S1, A2/S2, S1/S2, S1/S3, S2/S3, S3/F1, S1/F1, F1/F2, F1/P1, P1/P2}
    \draw[flecha, color=rTinta] (\a) -- (\b);
  \draw[flecha, color=rTinta] (S1.north) to[bend left=18] (F2.north west);
  \draw[flecha, color=rTinta] (F1) -- (P2);
  % leyenda
  \node[etiqueta, anchor=west, color=rGris] at (-1.25,-0.9) {Flecha: «la de destino necesita la de origen». Además, Toeplitz, Hankel y circulantes (A1.4--A1.6) se usan en S3, F1, P1 y P2.};
\end{tikzpicture}
\end{document}
